\section{Introduction}

Robots that take spoken commands in homes and competitions face the
cocktail-party problem with very little hardware. A common choice is the
ReSpeaker USB Mic Array v2.0: four omnidirectional microphones on a
32\,mm-radius circle with an on-chip DSP that outputs one processed channel
and a firmware direction-of-arrival (DOA) estimate. Our robot's speech
pipeline uses exactly that processed channel: it feeds a neural denoiser
(DeepFilterNet3 \cite{schroeter2023deepfilternet3}), then a voice activity
detector (VAD) built from an energy gate, an autocorrelation pitch test in
85--255\,Hz and an MFCC speaker lock, then speech recognition. Spatial
information is discarded, and the published firmware DOA is smoothed with
a linear average, which maps a talker crossing $0^\circ$ to $180^\circ$.

We ask what the raw microphones can add, in three steps a robot needs:
\emph{where} the talker is (angle), \emph{how far} (distance), and
\emph{attending} to them by suppressing other directions before the
denoiser. Our organising idea is that the VAD already provides the most
useful reference signal: frames without voice describe the room -- fans,
motors, diffuse noise -- and their cross-spectra can be used to clean the
localisation and to train the spatial filter.

Contributions: (i) a streaming localiser (pitch VAD $\to$ silence model
$\to$ SRP-PHAT $\to$ distance features $\to$ wrap-safe tracker) with an
ablation of each use of the silence model; (ii) an analysis of which
acoustic cues give distance on a 6.4\,cm array, showing that the
coherent-to-diffuse ratio works within a room and that a room covariate
(RT60) is what makes it transfer; (iii) a measurement of MVDR versus
delay-and-sum selectivity on this aperture with an estimated, not oracle,
steering direction; (iv) tools for the hardware study (live radar,
labelled recording, calibration) and an integration plan for a ROS~2
pipeline. All quantitative results here are from simulation; the hardware
evaluation is the next step and the paper says where we expect simulation
to be optimistic.
